\section{结果：接口为何需要来源治理}

结果分为两部分。首先，代表性记录说明测量侧能够把长时、多通道原始数据压缩为可回指的 observation packet，并以紧凑摘要形成明确的检索意图。其次，受控扩库结果表明，在相关补充资料加入后，语义近邻得分与核心权威来源的可见性可以发生背离；query 表述、语言、encoder 与分块方式会改变这种背离的幅度。本文据此讨论检索证据的可见性，而不将任何信号事件或候选解释提升为工程确认。

\subsection{代表性 trace：从多通道记录到受控检索意图}

以 A2 组的一条记录为例，初始化步骤生成了包含 597,607 个时间样本和 7 个中子电流通道的 active record；时间顺序校正后，记录覆盖约 7 d、名义采样间隔为 1 s，且未发现缺失值。确定性工具随后将事件、零值、通道关系和极值等结果写入记事板。该实例中，7 个通道各记录到 12 个孤立零值观测，短时偏离事件数为 30--34 个，通道平均值为 32.4 个。多个通道的零值时间重复出现，因而“同步零值—信号完整性/共享测量上下文”成为应保留在紧凑摘要中的检索线索，而不是可以直接命名的故障类型。

这一 trace 的关键不在于事件数本身，而在于每个摘要字段均可回指：零值时间、事件计数和持续时间来自事件扫描输出；通道间共变、极值位置和快照来自相应的确定性工具；condition report 只组织这些已保存字段。图~\ref{fig:trace-placeholder} 先以 V2 的事件计数汇总占位；正式版本将把一个局部时间—通道现象、相应的记事板字段、compact alert summary 和检索意图置于同一图中。该图不需要展示整段长记录，也不应将资料检索后的候选说明绘制成已确认根因。

\begin{figure*}[t]
  \centering
  \includegraphics[width=0.88\textwidth]{fig07_a2_event_summary}
  \caption{代表性 A2 记录的确定性事件汇总（V2 占位图）。柱形表示各通道在固定规则下得到的短时偏离事件与孤立零值计数，并非已确认故障计数。正式版本将以局部多通道时序、记事板字段和紧凑摘要替换。}
  \label{fig:trace-placeholder}
\end{figure*}

基于上述摘要，本案例的检索意图是查找与测量完整性、同步事件、冗余或共享仪表上下文有关的可定位技术段落。该意图将观测限定为可核查的技术主题；第五章后续的扩库实验进一步表明，仅以语义近邻选择这些段落并不能保证预设核心资料仍在结果中可见。


\subsection{字符分块扩库协议：语义相似度与核心来源可见性可以背离}

\texttt{nomic-embed-text} 字符分块扩库协议显示，补充资料的加入不必降低最高余弦相似度，却可能迅速降低核心来源在前 10 个结果中的占比。每个非零扩库规模在每个分块配置下均基于 100 次补充资料随机抽样取平均，随后再聚合 17 个已保存配置；下列 $\pm$ 值表示配置间的标准差。中文测量短语 query 的 CorePriority@10 从仅含核心资料时的 1.00，在加入 5 份补充资料后降至 $0.10\pm0.09$，加入 10 份后降至 0；与此同时，平均最高余弦相似度由 0.55 上升至 0.63。英文 query 的下降较缓：加入 5 份补充资料后 CorePriority@10 为 $0.86\pm0.05$，在加入 45 份后为 $0.38\pm0.05$，最高余弦相似度则维持在约 0.69--0.70。

这一结果不表示补充资料不相关。相反，它表明当部署要求核心权威来源保持可见时，相关补充资料本身就足以改变无来源约束语义检索的结果构成。图~\ref{fig:scope-priority-placeholder} 同时呈现最高余弦相似度与 CorePriority@10 随扩库规模的轨迹，使读者直接看到“相似度提高或稳定”与“核心来源可见性下降”可以同时发生。

\begin{figure*}[t]
  \centering
  \includegraphics[width=0.94\textwidth]{fig08_scope_priority}
  \caption{资料扩展下的核心来源可见性与最高余弦相似度（V2 占位图）。结果聚合 17 个已保存分块配置；阴影表示配置间标准差。核心来源可见性可以下降，即使最高余弦相似度上升或保持稳定。正式版本将统一为 V3 的中文术语与版式。}
  \label{fig:scope-priority-placeholder}
\end{figure*}

\subsection{MiniLM encoder 对照：query、语言与分块的交互}

在保持资料库、top-10 规则和来源指标不变的 MiniLM 核查中，以两种 encoder 和三种字符分块设置比较简短基线 query 与由 observation packet 线索形成的技术细节 query。在加入 10 份补充资料时，技术细节 query 相对基线 query 的 CorePriority@10 变化在英文设置中介于 $-0.03$ 至 $+0.14$，在中文设置中介于 $-0.23$ 至 $+0.02$。因此，向 query 中加入观察细节并非总能提高核心来源可见性。

在两种 MiniLM encoder 分别配合三种字符分块设置所形成的六种配置上，平均结果也体现了这种不对称性：英文技术细节 query 的 CorePriority@10 为 0.65，高于英文基线 query 的 0.60；中文技术细节 query 为 0.52，低于中文基线 query 的 0.62。完整扩库轨迹进一步显示，各曲线均从 core-only 条件下的 1.00 出发，但其下降速度、交叉位置和对 query 细节的响应随语言与配置而变。图~\ref{fig:configuration-placeholder} 以 V2 的热图占位呈现这一交互；正式版本将压缩为与 V3 叙述一致的分面热图，不再在正文同时保留多张语言重复图、梯度图和轨迹图。

\begin{figure*}[t]
  \centering
  \includegraphics[width=0.54\textwidth]{fig13_p0_priority_heatmap}
  \caption{不同 query 与 encoder—分块配置下的核心来源可见性（V2 占位图，$m=10$）。该图显示 query 表述的作用随语言与配置而改变。正式版本将减少 query 行数，并以中文标签呈现。}
  \label{fig:configuration-placeholder}
\end{figure*}

配置交互并不意味着无法作出选择。在本资料库和来源要求下，已记录结果能够筛出来源可见性保持较好的候选组合。表 5.1 列出两类具有不同用途的例子：\texttt{nomic-embed-text} 的 800/120 字符分块配合英文核心术语 query，在加入 10 份补充资料时的 CorePriority@10 为 $0.92\pm0.12$，在全部 51 份补充资料加入后仍为 0.60；多语 MiniLM 的 800/80 设置配合“高度同步且无时间滞后”的英文观测 query，在相应条件下为 $0.83\pm0.18$ 和 0.40。前者更适合作为保持核心来源可见性的检索基线，后者则表明来自 observation packet 的具体线索也能形成具有较高来源可见性的案例检索表述。

\begin{table*}[t]
\centering
\footnotesize
\begin{tabular}{p{0.16\textwidth} p{0.16\textwidth} p{0.22\textwidth} p{0.16\textwidth} p{0.16\textwidth}}
\toprule
候选用途 & encoder 与分块 & query 表述 & CorePriority@10（$m=10$） & CorePriority@10（$m=51$） \\
\midrule
核心来源可见性基线 & \texttt{nomic-embed-text}，800/120 字符 & \texttt{neutron} & $0.92\pm0.12$ & 0.60 \\
同步通道现象的案例检索 & 多语 MiniLM，800/80 字符 & \texttt{neutron current measurement highly synchronized channels without time lag} & $0.83\pm0.18$ & 0.40 \\
\bottomrule
\end{tabular}
\end{table*}


这些候选组合不是跨语料库、任务或评价目标的“最优配置”。特别是，短核心术语 query 的来源可见性较高，并不自动表示其最贴合某一条具体记录；而 observation-derived query 更接近案例现象，也仍需检验所返回段落的实际相关性。表 5.1 的作用是为检索协议提供可复核的条件性选择依据：在以核心来源可见性为目标时，应优先从这类在扩库后仍保持较高 CorePriority@10 的组合中选择，再结合段落相关性和工程复核确定实际部署配置。

MiniLM encoder 对照与 \texttt{nomic-embed-text} 字符分块扩库协议共同表明：在本研究固定的资料库、检索规则和已测试 encoder 范围内，相关补充资料扩展后核心来源可见性下降的风险并非单一表示模型所致。由于两类协议的 query 集和部分分块设置不同，本文不直接比较其绝对余弦分数，也不据此排列 encoder 优劣；比较重点是各自协议内 CorePriority@10 随扩库规模的变化。技术细节 query 在不同语言和配置下的作用并不一致，因此结构化摘要应被理解为受控的检索表述，而非必然提高权威来源可见性的查询扩展。

\subsection{分块边界与检索单元核查}

固定词数分块的 MiniLM 对照表明，上述风险并非只出现在 \texttt{nomic-embed-text} 的字符分块中。多语种 MiniLM 对中文短语产生的最高余弦相似度约为 0.62--0.64，高于另一 MiniLM 的约 0.27；但随着补充资料加入，两种 encoder 的 CorePriority@10 均出现下降。语义表示对不同语言的近邻得分具有实质影响，却不能替代来源要求。

\texttt{nomic-embed-text} 索引上的 self-retrieval 则说明，检索单元的定义也会影响“找回”的含义。在一个 800/50 设置的 core-only 条件下，exact-chunk Recall@10 为 0.7625，而 same-document Recall@10 为 0.8625。该差异来自重叠分块：原始段落未必以同一 node 返回，但同一文档中的相邻段落可能已包含对应信息。因此，self-retrieval 结果用于解释 chunk 边界与段落定位，而不应与 CorePriority@10 混为同一种证据。

\begin{figure*}[t]
  \centering
  \includegraphics[width=0.84\textwidth]{fig17_self_retrieval_recall}
  \caption{self-retrieval 的 Recall@10 结果（V2 占位图）。该核查区分 exact-chunk 与 same-document 两种匹配标准，用于解释重叠分块下检索单元与段落定位的关系，而不衡量工程正确性。正式版本可与 MRR 合并重绘，或移至补充材料。}
  \label{fig:self-retrieval-placeholder}
\end{figure*}

\subsection{小结}

代表性 trace 说明，测量侧能够将无标签记录转换为可回指的检索线索；\texttt{nomic-embed-text} 字符分块扩库协议说明，证据侧若只依赖语义近邻，核心权威来源的可见性可能随相关资料的加入而丢失。MiniLM encoder 对照进一步表明：在本研究已测试的 encoder 范围内，这一风险并非单一向量表示所致；其幅度仍受 query、语言和分块方式共同影响，不能化约为某个组件的单一最优选择。与此同时，扩库后的来源可见性指标能够筛出条件性候选配置，为检索协议选择提供实证依据。由此，来源 metadata 与显式来源策略不是检索实现的附属信息，而是 measurement-to-evidence interface 中需要被设计和评价的部分。
